\documentclass[aspectratio=169,11pt]{beamer}

% ---------------------------------------------------------------
% Encoding und Sprache (xelatex required)
% ---------------------------------------------------------------
\usepackage{fontspec}
\usepackage[ngerman,provide=*]{babel}
\usepackage{csquotes}

% ---------------------------------------------------------------
% Polices : Inconsolata (mono) pour titres, Open Sans pour le corps
% ---------------------------------------------------------------
\setmainfont{Open Sans}[
  Scale=0.96,
  BoldFont={Open Sans Bold},
  ItalicFont={Open Sans Italic},
  BoldItalicFont={Open Sans Bold Italic},
]
\setmonofont{Inconsolata}[Scale=0.95]
% \titlemono pour les titres en mono.
% \mdash : en-dash en DejaVu Sans Mono (Inconsolata Medium n'a pas U+2013).
\newfontfamily\titlemono{Inconsolata}[Scale=0.96]
\newfontfamily\monodash{DejaVu Sans Mono}[Scale=0.95]
\newcommand{\mdash}{\,{\monodash–}\,}

% ---------------------------------------------------------------
% Palette Solarized dark (Ethan Schoonover)
% ---------------------------------------------------------------
\usepackage{xcolor}
\definecolor{solBase03}{HTML}{002b36}
\definecolor{solBase02}{HTML}{073642}
\definecolor{solBase01}{HTML}{586e75}
\definecolor{solBase00}{HTML}{657b83}
\definecolor{solBase0}{HTML}{839496}
\definecolor{solBase1}{HTML}{93a1a1}
\definecolor{solBase2}{HTML}{eee8d5}
\definecolor{solBase3}{HTML}{fdf6e3}
\definecolor{solYellow}{HTML}{b58900}
\definecolor{solOrange}{HTML}{cb4b16}
\definecolor{solRed}{HTML}{dc322f}
\definecolor{solMagenta}{HTML}{d33682}
\definecolor{solViolet}{HTML}{6c71c4}
\definecolor{solBlue}{HTML}{268bd2}
\definecolor{solCyan}{HTML}{2aa198}
\definecolor{solGreen}{HTML}{859900}

% ---------------------------------------------------------------
% Beamer theme: Solarized dark
% ---------------------------------------------------------------
\usetheme{default}
\setbeamertemplate{navigation symbols}{}

% Fond et texte par défaut
\setbeamercolor{background canvas}{bg=solBase03}
\setbeamercolor{normal text}{fg=solBase0,bg=solBase03}
\setbeamercolor{structure}{fg=solCyan}
\setbeamercolor{alerted text}{fg=solOrange}
\setbeamercolor{example text}{fg=solGreen}

% Accents renforcés : \emph en cyan italique, \textbf en blanc cassé
\let\oldemph\emph
\renewcommand{\emph}[1]{{\color{solCyan}\oldemph{#1}}}
\let\oldtextbf\textbf
\renewcommand{\textbf}[1]{{\color{solBase3}\oldtextbf{#1}}}

% Titres et sous-titres de slide
\setbeamerfont{frametitle}{family=\titlemono,size=\large,series=\bfseries}
\setbeamercolor{frametitle}{fg=solBase3,bg=solBase03}
\setbeamerfont{title}{family=\titlemono,size=\Large,series=\bfseries}
\setbeamercolor{title}{fg=solBase3}
\setbeamerfont{subtitle}{family=\titlemono,size=\normalsize}
\setbeamercolor{subtitle}{fg=solBase1}
\setbeamerfont{author}{size=\small}
\setbeamercolor{author}{fg=solBase1}
\setbeamercolor{institute}{fg=solBase01}
\setbeamercolor{date}{fg=solBase01}

% Bullets
\setbeamertemplate{itemize item}{\color{solYellow}\rule[0.4ex]{0.5ex}{0.5ex}}
\setbeamertemplate{itemize subitem}{\color{solCyan}\textendash}
\setbeamertemplate{enumerate item}{\color{solYellow}\bfseries\insertenumlabel.}

% Blocks
\setbeamercolor{block title}{fg=solBase3,bg=solBase02}
\setbeamercolor{block body}{fg=solBase0,bg=solBase02}
\setbeamerfont{block title}{family=\titlemono,series=\bfseries,size=\normalsize}

% Footer discret
\setbeamertemplate{footline}{%
  \leavevmode%
  \hbox{%
    \begin{beamercolorbox}[wd=.55\paperwidth,ht=2.25ex,dp=1ex,leftskip=1em]{footer}%
      \footnotesize\color{solBase01}\titlemono\small MTA\mdash{}Sitzung 1\normalsize%
    \end{beamercolorbox}%
    \begin{beamercolorbox}[wd=.45\paperwidth,ht=2.25ex,dp=1ex,rightskip=1em,leftskip=2ex]{footer}%
      \hfill\footnotesize\color{solBase01}\titlemono\insertframenumber\,/\,\inserttotalframenumber%
    \end{beamercolorbox}}%
}

% ---------------------------------------------------------------
% Autres paquets
% ---------------------------------------------------------------
\usepackage{booktabs}
\usepackage{array}
\usepackage{tabularx}
\usepackage{colortbl}
\usepackage{amsmath}
\usepackage{tikz}
\usetikzlibrary{positioning,arrows.meta,calc,matrix,fit,backgrounds}
\arrayrulecolor{solBase01}

% ---------------------------------------------------------------
% Titelinformationen
% ---------------------------------------------------------------
\title{Topic-Modelle als qualitative Methode?}
\subtitle{Sitzung 1\mdash{}Positionierung gegenüber\\objektiver Hermeneutik und qualitativer Inhaltsanalyse}
\author{Prof.~Dr.~Christian Papilloud}
\institute{Institut für Soziologie\\Martin-Luther-Universität Halle-Wittenberg}
\date{Master Soziologie – 2026}

\begin{document}

% ===============================================================
\begin{frame}[plain]
  \titlepage
\end{frame}

% ===============================================================
\begin{frame}{Worum es in dieser Sitzung geht}

  Bevor wir zur praktischen Anwendung des Topic-Modell-Ansatzes mit
  dem Programm MTA (Multi-Text Analyser) kommen, beschäftigen wir
  uns in drei Sitzungen mit dem \emph{theoretischen Hintergrund}
  von Topic-Modellen.
  \vfill
  Drei Sitzungen, drei Fragen:
  \begin{enumerate}
    \item[\textbf{1.}] \textbf{Heute:} Ist Topic-Modellierung eine qualitative
          Methode? Wenn ja, in welchem Sinn? Wie verhält sie sich zu klassischen
          qualitativen Verfahren der Soziologie?
    \item[\textbf{2.}] \textbf{Nächste Sitzung:} Wie funktionieren Topic-Modelle
          \emph{mathematisch}? Was ist heute Stand der Technik (BERT, LLMs)?
    \item[\textbf{3.}] \textbf{Dritte Sitzung:} Was können Topic-Modelle nicht,
          und welche Konsequenzen folgen daraus für unsere Forschungspraxis?
  \end{enumerate}

\end{frame}

% ===============================================================
\begin{frame}{Eine erste, kontraintuitive These}

  \begin{block}{Papilloud \& Hinneburg (2018, S.~1)}
    \enquote{Topic-Modelle stellen ein quantitativ basiertes Werkzeug für
    die Textanalyse dar [\ldots]. Topic-Modell-Verfahren sind zwar quantitativ
    basiert, aber sie \emph{dienen der qualitativen Forschung} – eine Aussage,
    die nicht sofort einleuchten mag.}
  \end{block}

  \vfill
  Warum nicht? Weil das Begriffspaar \enquote{quantitativ vs.\ qualitativ}
  meistens als Gegensatz verstanden wird:
  \begin{itemize}
    \item Statistik $\leftrightarrow$ Sinnverstehen
    \item Zahlen $\leftrightarrow$ Texte
    \item Erklären $\leftrightarrow$ Verstehen
  \end{itemize}

  \vfill
  Ein Verfahren, das Häufigkeiten zählt und Matrizen faktorisiert,
  klingt nicht nach hermeneutischer Arbeit. Und doch \ldots\

\end{frame}

% ===============================================================
\begin{frame}{Drei gemeinsame Grundannahmen}

  Objektive Hermeneutik, qualitative Inhaltsanalyse \emph{und}
  Topic-Modell-Verfahren teilen drei Voraussetzungen:

  \vfill
  \begin{enumerate}
    \item Begriffe und Sätze werden \textbf{nicht zufällig} verwendet.
          Sinn ist nicht zufällig.
    \vfill
    \item Texte enthalten \textbf{Strukturen}, die ihrem Sinn zugrunde
          liegen. Aufgabe der Analyse: diese Strukturen rekonstruieren.
    \vfill
    \item Die \textbf{Deutung} – und damit der/die Forschende – spielt
          eine Rolle, sowohl bei der Organisation des Verfahrens als auch
          bei der Interpretation der Ergebnisse.
  \end{enumerate}

  \vfill
  \footnotesize\color{black!60}
  Vgl. Papilloud \& Hinneburg 2018, S.~2–3.

\end{frame}

% ===============================================================
\begin{frame}{Aber: drei Methoden, drei Antworten}

  Die gleichen Annahmen führen zu sehr unterschiedlichen Antworten
  auf drei Fragen:

  \vfill
  \begin{itemize}
    \item Was heißt \enquote{Struktur} des Sinns? \\
          $\rightarrow$ \textit{latente Sinnstruktur} (Oevermann) vs.\
          \textit{Kategoriensystem} (Mayring) vs.\
          \textit{topics} (Blei, Hinneburg)
    \vfill
    \item Wie wird \enquote{Sinn} algorithmisch erzeugt? \\
          $\rightarrow$ universelle Grammatik (Chomsky) vs.\
          Worthäufigkeiten und Konfigurationen
    \vfill
    \item Welche Rolle spielt der Forschende? \\
          $\rightarrow$ Gedankenexperiment vs.\
          Codierung vs.\ Interpretation eines Klassifikationsmusters
  \end{itemize}

\end{frame}

% ===============================================================
\section{1. Objektive Hermeneutik (Oevermann)}

\begin{frame}{Was uns interessiert: nicht die ganze Theorie}

  Die objektive Hermeneutik ist ein elaboriertes Forschungsprogramm.
  Wir konzentrieren uns auf drei Punkte, die für den \emph{Kontrast}
  mit Topic-Modellen entscheidend sind:

  \vfill
  \begin{enumerate}
    \item Was ist \textbf{Lebenspraxis} und \textbf{objektive Sinnstruktur}?
    \vfill
    \item Welches \textbf{methodologische Verfahren} folgt daraus?
    \vfill
    \item Welche \textbf{Vorstellung von Sprache und Sinn} liegt zugrunde?
  \end{enumerate}

  \vfill
  \footnotesize\color{black!60}
  Quelle: Oevermann 1995, 1996, 2000, 2002, 2008, 2010.
  Vgl.\ Papilloud \& Hinneburg 2018, Abschnitt 1.1.

\end{frame}

% ===============================================================
\begin{frame}{Lebenspraxis und objektive Sinnstruktur}

  Bei Oevermann ist \textbf{Lebenspraxis} das Zusammenspiel
  physischer, psychischer, sozialer und kultureller Potentiale,
  die einem Einzelakteur die Gestaltung seines Lebens erlauben.

  \vfill
  \begin{block}{Zentrale Idee}
    Jede Lebenspraxis ist eine \textbf{Sequenz von Entscheidungen} –
    nicht eine Sequenz von Aussagen.
  \end{block}

  \vfill
  \begin{itemize}
    \item Der Akteur steht ständig vor \textbf{Krisen}:
          Entscheidungssituationen ohne Richtig-/Falsch-Kalkulation.
    \item Er muss sich entscheiden – \enquote{er kann sich dem Prinzip
          nach nicht nicht entscheiden}.
    \item Die Begründung kommt \emph{retrospektiv}:
          \enquote{widersprüchliche Einheit von Entscheidungszwang
          und Begründungsverpflichtung}.
    \item Die Verkettung dieser Entscheidungen bildet eine
          \textbf{objektive Sinnstruktur}.
  \end{itemize}

\end{frame}

% ===============================================================
\begin{frame}{Manifest und latent}

  \begin{columns}[T,onlytextwidth]
    \column{0.48\textwidth}
    \textbf{Manifeste Sinnstruktur} \\[2pt]
    Was der Akteur \emph{explizit} sagt. Was im Protokoll dokumentiert
    ist. Beobachtbar.

    \vspace{1.5em}
    \textbf{Latente Sinnstruktur} \\[2pt]
    Die \emph{Verkettung} der Aussagen, das, was sich aus ihr
    \emph{ergibt}. Nicht direkt beobachtbar. Muss vom Forschenden
    \emph{gedeutet} werden.

    \column{0.48\textwidth}
    \begin{block}{Pointe}
      Die latente Sinnstruktur ist nicht das, was der Akteur
      \emph{wirklich gemeint} hat. Sie ist die Struktur, die
      seinen Aussagen \textbf{vorausgeht} und sie ermöglicht –
      \enquote{regelerzeugter objektiver Sinn} (Oevermann 2002).
    \end{block}

    \vspace{0.5em}
    \footnotesize\color{black!60}
    Annahme: Es gibt eine \emph{universelle generative Grammatik}
    (Chomsky), die Sinn algorithmisch erzeugt.
  \end{columns}

\end{frame}

% ===============================================================
\begin{frame}{Das Verfahren: Sequenzanalyse}

  \textbf{Schritte der Fallrekonstruktion:}

  \vfill
  \begin{enumerate}
    \item Auswahl eines Protokolls (Transkript, Beobachtungs\-protokoll
          usw.) bzw.\ eines \emph{Falls}.
    \vfill
    \item \textbf{Sequenzanalyse}: das Protokoll wird Schritt für
          Schritt durchgegangen. An jeder Sequenzstelle werden
          \emph{Lesarten} entwickelt – alle möglichen Deutungen,
          die in diesem Kontext denkbar wären.
    \vfill
    \item Schrittweise Eliminierung: nur die Lesarten überleben,
          die mit dem weiteren Verlauf des Protokolls vereinbar sind.
    \vfill
    \item Ziel: die \textbf{Fallstrukturgesetzlichkeit} des
          beobachteten Falls rekonstruieren.
  \end{enumerate}

  \vfill
  \footnotesize\color{black!60}
  Ein arbeitsintensives Verfahren: stundenlange Interpretation
  einzelner Sätze, oft in Forschungsgruppen.

\end{frame}

% ===============================================================
\section{2. Qualitative Inhaltsanalyse (Mayring)}

\begin{frame}{Mayrings Programm: ein \emph{hybrides} Profil}

  Die qualitative Inhaltsanalyse (Mayring, seit Anfang der 1980er)
  möchte \emph{nicht} im Gegensatz zur Quantifizierung stehen,
  sondern beide verbinden:

  \vfill
  \begin{block}{Mayring 2002, S.~114}
    \enquote{Texte systematisch analysieren, indem sie das Material
    schrittweise mit theoriegeleitet am Material entwickelten
    Kategoriensystemen bearbeitet.}
  \end{block}

  \vfill
  Kernoperation: das \textbf{Kategoriensystem}. \\[4pt]
  Bestandteile (Mayring 1994):
  \begin{itemize}
    \item Kategorien und Unterkategorien
    \item Kategoriendefinitionen
    \item Ankerbeispiele
    \item begründete Überarbeitung im Laufe der Auswertung
  \end{itemize}

\end{frame}

% ===============================================================
\begin{frame}{Was das Kategoriensystem leisten soll}

  Zwei Ansprüche, die sich \emph{spannungsreich} ergänzen:

  \vfill
  \begin{enumerate}
    \item \textbf{Hermeneutische Tradition:}
          das Kategoriensystem soll den \emph{manifesten} \emph{und}
          den \emph{latenten} Sinn der Textquelle abbilden.
    \vfill
    \item \textbf{Quantitative Tradition:}
          das Kategoriensystem soll eine \emph{systematische
          Überprüfung} von Strukturen erlauben (Reliabilität,
          Interkoderreliabilität, Triangulation).
  \end{enumerate}

  \vfill
  \begin{block}{Ziel}
    \enquote{Intersubjektiv nachvollziehbare Ergebnisse}: das, was
    eine Codiererin findet, soll auch eine zweite, unabhängige
    Codiererin finden können – Transparenz, Vergleichbarkeit.
  \end{block}

\end{frame}

% ===============================================================
\begin{frame}{Bekannte Kritik an Mayring}

  Die qualitative Inhaltsanalyse wird häufig kritisiert.
  Drei Punkte aus dem Buch (Papilloud \& Hinneburg 2018, S.~11–12):

  \vfill
  \begin{itemize}
    \item Die Einbindung statistischer Methoden bleibt
          \textbf{rudimentär}: meistens auf Auszählung von
          Kategorienhäufigkeiten beschränkt.
    \vfill
    \item Kategorien werden zwar \emph{induktiv} aus dem Text gewonnen,
          sind aber zunächst \emph{deduktiv} angesetzt – der Bezug
          zum Text als Ganzem geht verloren (Plant 1996,
          Froggatt 2001, Lamnek \& Krell 2005).
    \vfill
    \item Nach Kruse (2015): teilweise eine \enquote{rein quantifizierende
          Subsumtionslogik}. Mayring würde Äußerungen
          \emph{inventarisieren}, aber nicht die Ebene der
          \emph{latenten} Strukturen erreichen.
  \end{itemize}

\end{frame}

% ===============================================================
\section{3. Was Topic-Modelle anders machen}

\begin{frame}{Wo Topic-Modelle stehen}

  Topic-Modelle teilen mit \emph{beiden} Methoden die drei
  Grundannahmen vom Anfang. Aber sie unterscheiden sich von
  beiden in einem entscheidenden Punkt:

  \vfill
  \begin{block}{Was sich grundsätzlich ändert}
    Bei Oevermann und Mayring ist das \emph{Kategoriensystem}
    (oder die Fallstruktur) das \textbf{Ergebnis menschlicher
    Interpretation}.
    \\[4pt]
    Bei Topic-Modellen ist es das \textbf{Ergebnis eines mathematischen
    Algorithmus} – der dann seinerseits interpretiert werden muss.
  \end{block}

  \vfill
  Was bedeutet das konkret?

\end{frame}

% ===============================================================
\begin{frame}{Drei Unterschiede zur objektiven Hermeneutik}

  \begin{enumerate}
    \item \textbf{Begriff von \enquote{Algorithmus}.} \\
          Bei Oevermann: regelgeleitete Selektion in
          Möglichkeitswelten, universelle generative Grammatik.\\
          Bei TM: statistische / linearealgebraische Verfahren –
          \emph{die Grammatik ist nicht universell}, sondern verändert
          sich mit Gebrauch und Kontext.
    \vfill
    \item \textbf{Was wird gesucht?} \\
          Bei Oevermann: die \emph{Fallstrukturgesetzlichkeit} eines
          bestimmten Falls.\\
          Bei TM: \emph{Klassifikationsmuster} über viele Dokumente
          hinweg.
    \vfill
    \item \textbf{Rolle der Kontextvariation.} \\
          Bei Oevermann: Gedankenexperimente, Lesarten – viele
          alternative Kontexte werden \emph{aktiv} mobilisiert.\\
          Bei TM: solche Variationen werden \emph{bewusst vermieden}
          (Sinclair 1992): es zählen die typischen Merkmale des Textes.
  \end{enumerate}

\end{frame}

% ===============================================================
\begin{frame}{Nähe zur qualitativen Inhaltsanalyse}

  Topic-Modelle stehen Mayring \emph{näher} als Oevermann.
  Beide:
  \vfill
  \begin{itemize}
    \item bauen \textbf{Kategoriensysteme} aus dem Material auf,
    \item streben \textbf{Intersubjektivität} an (zweimal angewandt,
          zweimal vergleichbares Ergebnis),
    \item nutzen \textbf{Häufigkeiten} – aber für TM ist das nicht
          ein nachträglicher quantitativer Schritt, sondern die
          \emph{Grundlage} der Kategorienbildung.
  \end{itemize}

  \vfill
  \begin{block}{Wichtige Differenz}
    Was Mayring kritisiert wird (Kategorien bleiben oberflächlich,
    erreichen die latenten Strukturen nicht), versuchen TM
    methodisch zu vermeiden:
    \\[4pt]
    Sie nehmen die \textbf{Indexikalität der Sprache} ernst.
    Bedeutung ergibt sich aus der \emph{Konfiguration} von Begriffen,
    nicht aus ihrer isolierten Codierung.
  \end{block}

\end{frame}

% ===============================================================
\begin{frame}{Synopse}

  \footnotesize
  \begin{tabularx}{\textwidth}{@{}p{0.20\textwidth}XXX@{}}
    \toprule
    & \textbf{Objektive Hermeneutik} & \textbf{Qual.\ Inhaltsanalyse} & \textbf{Topic-Modell} \\
    \midrule
    \textbf{Gegenstand}
      & Einzelfall, biographische Sequenz
      & Textmaterial, Inventar von Aussagen
      & Textsammlung, Korpus \\[6pt]
    \textbf{Algorithmus}
      & generative Grammatik (Chomsky)
      & Kategoriensystem, manuelle Codierung
      & statistische / lineare Algebra \\[6pt]
    \textbf{Sinnbegriff}
      & latent, regelerzeugt, objektiv
      & manifest \emph{und} latent
      & latent als Konfiguration von Begriffen \\[6pt]
    \textbf{Forschende:r}
      & deutet, entwickelt Lesarten
      & codiert anhand des Systems
      & interpretiert das Klassifikationsmuster \\[6pt]
    \textbf{Ergebnis}
      & Fallstruktur\-gesetzlichkeit
      & Kategorien mit Häufigkeiten
      & topics + Dokument\-zugehörigkeiten \\[6pt]
    \textbf{Korpusgröße}
      & klein (1–20 Fälle)
      & mittel
      & groß (Hunderte bis Millionen) \\
    \bottomrule
  \end{tabularx}

\end{frame}

% ===============================================================
\section{Zwischenfazit und Ausblick}

\begin{frame}{Zwischenfazit der heutigen Sitzung}

  \begin{enumerate}
    \item Die Frage \enquote{ist TM qualitativ?} hat eine differenzierte
          Antwort: \emph{ja}, im Sinne der drei Grundannahmen –
          \emph{aber} mit einem mathematischen Algorithmus an der Stelle,
          wo bei klassischen Methoden ein interpretativer Akt steht.
    \vfill
    \item TM ist \textbf{näher an Mayring als an Oevermann}.
          Es teilt das Ziel der Intersubjektivität und der
          Kategoriensysteme, geht aber methodisch anders vor.
    \vfill
    \item Der \emph{interpretative Akt} verschwindet \emph{nicht} –
          er verschiebt sich nur. Bei TM geht es darum, das
          \textbf{Klassifikationsmuster} zu deuten, das der Algorithmus
          erzeugt hat. Dazu kommen wir in der dritten Sitzung zurück.
  \end{enumerate}

\end{frame}

% ===============================================================
\begin{frame}{Was uns nächste Woche erwartet}

  \textbf{Sitzung 2 – Wie funktionieren Topic-Modelle?}

  \vfill
  \begin{itemize}
    \item Was ist ein \emph{topic} mathematisch? Die Idee der
          Matrixfaktorisierung anhand eines kleinen Beispiels.
    \vfill
    \item Vergleich mit Clusteranalyse und Faktorenanalyse.
    \vfill
    \item Die zwei Hauptverfahren: \textbf{LDA} (probabilistisch,
          bayessch) und \textbf{NMF} (lineare Algebra).
    \vfill
    \item \emph{Neue Welt:} \textbf{BERTopic} und der Sprung zu
          \emph{Embeddings}. Was bringen LLMs (GPT, Claude) in die
          Diskussion ein?
    \vfill
    \item Vorgriff auf die Frage, die uns durch die Praxis mit MTA
          begleiten wird: \emph{warum} verwenden wir hier (noch)
          NMF und LDA – und nicht BERTopic oder einen LLM?
  \end{itemize}

\end{frame}

% ===============================================================
\begin{frame}{Eine Frage, die wir offen lassen}

  \begin{block}{Zum Mitnehmen}
    Eine Sozialwissenschaftlerin, die mit MTA arbeitet, codiert nicht
    selbst, sondern sie \emph{liest die Codierung} einer Maschine.
    \\[4pt]
    Verschiebt das die Subjektivität der Interpretation
    – oder verlagert es sie nur an einen anderen, weniger sichtbaren
    Ort?
  \end{block}

  \vfill
  Diese Frage werden wir in der dritten Sitzung wieder aufnehmen,
  wenn wir den Schlussblick des Buches (Kap.~8: \emph{Rück- und
  Ausblick}) besprechen.

\end{frame}

% ===============================================================
\begin{frame}[allowframebreaks]{Zitierte Literatur (Auswahl)}

  \footnotesize
  \begin{itemize}
    \item \textbf{Papilloud, C., \& Hinneburg, A.} (2018).
          \emph{Einführung in die qualitative Analyse von Texten mit
          Topic-Modellen}. Wiesbaden: Springer.
          \\\textit{Grundlage dieser Sitzung; bes.\ Kap.~1.}
    \item Oevermann, U. (1995). Ein Modell der Struktur von Religiosität.
          In: M.\ Wohlrab-Sahr (Hrsg.), \emph{Biographie und Religion},
          27–102. Frankfurt: Campus.
    \item Oevermann, U. (2002). \emph{Klinische Soziologie auf der Basis
          der Methodologie objektiver Hermeneutik – Manifest}.
    \item Oevermann, U. (2008). \emph{Krise und Routine als analytisches
          Paradigma in den Sozialwissenschaften} (Abschiedsvorlesung).
    \item Mayring, P. (1994). Qualitative Inhaltsanalyse. In: A.\ Boehm
          et al.\ (Hrsg.), \emph{Texte verstehen}, 159–175. Konstanz: UVK.
    \item Mayring, P. (2002). \emph{Einführung in die qualitative
          Sozialforschung}. Basel: Beltz.
    \item Kruse, J. (2015). \emph{Qualitative Interviewforschung. Ein
          integrativer Ansatz}. Weinheim: Beltz.
    \item Przyborski, A., \& Wohlrab-Sahr, M. (2009). \emph{Qualitative
          Sozialforschung. Ein Arbeitsbuch}. München: Oldenbourg.
    \item Chomsky, N. (1981). \emph{Regeln und Repräsentationen}.
          Frankfurt: Suhrkamp.
    \item Sinclair, J. (1992). The Automatic Analysis of Corpora.
          In: J.\ Svartvik (Hrsg.), \emph{Directions in Corpus Linguistics}.
  \end{itemize}

\end{frame}

% ===============================================================
\begin{frame}[plain]
  \begin{center}
    \vspace*{2em}
    {\Large\bfseries Vielen Dank.}
    \\[2em]
    Fragen, Einwände, Anschlüsse?
    \\[3em]
    \footnotesize\color{black!60}
    christian.papilloud@soziologie.uni-halle.de
  \end{center}
\end{frame}

\end{document}
